\subsection{Ergodic unitarily invariant measures on infinite Hermitian matrices}\label{section1.2}

Let $\mathbb{U}(N)$ and $H(N)$ be the group of $N \times N$ unitary matrices and the space of $N \times N$ Hermitian matrices respectively. Let $\mathbb{U}(\infty)$ be the inductive limit of unitary groups
\begin{gather*}
\mathbb{U}(\infty)=\mathop{\underset{\rightarrow}{\lim}}\mathbb{U}(N),
\end{gather*}
 under the natural embeddings. In more explicit terms an element of $\mathbb{U}(\infty)$ is an infinite block matrix whose top corner is an $N\times N$ unitary matrix for some finite $N$ and the other block is the (infinite) identity matrix.

Consider the so called corners maps $\pi^N_{N-1}\colon H(N)\to H(N-1)$ defined by
\begin{align}\label{Corners}
\pi_{N-1}^N\big[(h_{ij})_{i,j=1}^N\big]= (h_{ij} )_{i,j=1}^{N-1}.
\end{align}
Let $H$ be the space of infinite Hermitian matrices defined as the projective limit
\begin{gather*}
H=\mathop{\underset{\leftarrow}{\lim}}H(N),
\end{gather*}
under the corners maps. Moreover, let $H_+(N)\subset H(N)$ denote the space of $N \times N$ positive-definite matrices; namely matrices with positive eigenvalues. By Cauchy's interlacing theorem we get that $\pi^N_{N-1}\colon H_+(N)\to H_+(N-1)$. Thus, we can also correctly define the projective limit $H_+=\underset{\leftarrow}{\lim}H_+(N) \subset H$.

Now, $\mathbb{U}(\infty)$ acts on $H$ by conjugation: for each $u \in \mathbb{U}(\infty)$ we have a map $T_u\colon H \to H$ given by $T_u(h)=u^*hu$. It is a classical theorem of Pickrell \cite{Pickrell} and also Olshanski and Vershik \cite{OlshanskiVershik} that ergodic measures on $H$ for this action (namely ones such that all $\mathbb{U}(\infty)$-invariant subsets have mass 0 or 1) are parametrized by the infinite dimensional space $\Omega \subset \mathbb{R}^{2\infty+2}$:
\begin{gather}
\Omega =\Big\{\omega=(\alpha^+,\alpha^-,\gamma_1,\delta)\in \mathbb{R}^{2\infty+2}=\mathbb{R}^\infty \times \mathbb{R}^\infty \times \mathbb{R} \times \mathbb{R}\,|\nonumber\\
\hphantom{\Omega =\Big\{}{}\alpha^+ =(\alpha_1^+\ge \alpha_2^+\ge \cdots \ge 0) ; \, \alpha^-=(\alpha_1^-\ge \alpha_2^-\ge \cdots \ge 0);\nonumber\\
\hphantom{\Omega =\Big\{}{} \gamma_1 \in \mathbb{R} ;\, \delta \ge 0 ;\,\sum_{}^{}(\alpha_i^+)^2 + \sum_{}^{}(\alpha_i^-)^2\le \delta \Big\},\nonumber\\
 \gamma_2=\delta-\sum_{}^{}\big(\alpha_i^+\big)^2 - \sum_{}^{}\big(\alpha_i^-\big)^2.\label{Omegadefinition}
\end{gather}
Observe that $\Omega$ is a locally compact space. We then have the following classification theorem, see \cite{OlshanskiVershik,Pickrell}, also \cite{Defosseux}.

\begin{Theorem}[Pickrell, Olshanski--Vershik] There exists a parametrization of ergodic $\mathbb{U}(\infty)$-invariant probability measures on the space $H$ by the points of the space $\Omega$ described as follows. Given $\omega \in \Omega$ the characteristic function of the ergodic measure $M_{\omega}$ is given by
\begin{gather*}
\int_{X \in H}^{}{\rm e}^{{\rm i}\operatorname{Tr}(\operatorname{diag}(r_1,\dots,r_n,0,0,\dots)X)}M_{\omega}({\rm d}X)=\prod_{j=1}^{n}F_{\omega}(r_j),
\end{gather*}
where
\begin{gather*}
F_{\omega}(x)={\rm e}^{{\rm i}\gamma_1x-\frac{\gamma_2}{2}x^ 2}\prod_{k=1}^{\infty}\frac{{\rm e}^{-{\rm i}\alpha_k^+x}}{1-{\rm i}\alpha_k^+x}\prod_{k=1}^{\infty}\frac{{\rm e}^{{\rm i}\alpha_k^-x}}{1+{\rm i}\alpha_k^-x} .
\end{gather*}
\end{Theorem}
Note that the characteristic function $F_{\omega}$ is well defined for all $\omega \in \Omega$ by the fact that the sum of squares of the $\alpha$'s is finite. Also, observe that parameter $\gamma_2$ corresponds to an infinite random Hermitian matrix with independent (subject to the Hermitian constraint) complex Gaussian entries with mean 0 and variance $\gamma_2$ on the diagonal.

\subsection[The inverse Wishart measures $\mathsf{M}^{(\nu)}$ on infinite positive-definite matrices]{The inverse Wishart measures $\boldsymbol{\mathsf{M}^{(\nu)}}$ on infinite positive-definite matrices}\label{section1.3}

For $\nu>-1$, consider the complex Wishart (or Laguerre ensemble) probability measure on $N \times N$ Hermitian matrices, supported on $H_+(N)$:
\begin{gather*}
\mathcal{M}^{(\nu),N}({\rm d}Y)=\widetilde{\operatorname{const}_{\nu, N}}\det(Y)^{\nu} {\rm e}^{-\operatorname{Tr}Y}\mathbf{1}_{\{Y \in H_+(N)\}}{\rm d}Y\,
\end{gather*}
where throughout the paper for a Hermitian matrix $Y$, ${\rm d}Y$ denotes Lebesgue measure on $H(N)$ (we suppress dependence on~$N$):
\begin{gather*}
{\rm d}Y=\prod_{j=1}^{N}{\rm d}Y_{jj} \prod_{1 \le j <k \le N}^{}{\rm d} \operatorname{Re} Y_{jk}{\rm d} \operatorname{Im} Y_{jk}.
\end{gather*}
The restriction $\nu>-1$ is so that the normalization constant $\widetilde{\operatorname{const}_{\nu, N}}$ is finite.

Under the change of variables $Y=2X^{-1}$ we obtain the central object of study in this paper, the inverse Wishart probability measures on $H_+(N)$
\begin{align}\label{MeasureDefinition}
\mathsf{M}^{(\nu),N}({\rm d}X)= \operatorname{const}_{\nu, N} \det(X)^{-\nu-2N} {\rm e}^{-2\operatorname{Tr}X^{-1}}\mathbf{1}_{ \{X \in H_+(N) \}}{\rm d}X.
\end{align}
Observe that, for all $N \in \mathbb{N}$ the measure $\mathsf{M}^{(\nu),N}$ is unitarily invariant, namely invariant under the action of $\mathbb{U}(N)$ by conjugation. Then, we have the following consistency result:

\begin{Proposition}\label{ConsistencyIntro} For $\nu >-1$ the measures $\big\{\mathsf{M}^{(\nu),N}\big\}_{N\ge 1}$ form a projective family
\begin{gather*}
\big(\pi_{N-1}^N\big)_*\mathsf{M}^{(\nu),N}=\mathsf{M}^{(\nu),N-1}.
\end{gather*}
\end{Proposition}

Thus, by Kolmogorov's theorem we obtain a unitarily invariant measure $\mathsf{M}^{(\nu)}$ on $H$ that is supported on $H_{+}$.

Proposition \ref{ConsistencyIntro} is proven in Section~\ref{section2} as Proposition \ref{consistencymatrices}. A key role in the proof is played by the Bessel orthogonal polynomials \cite{BesselPolynomials}.

\begin{Remark}Although we have not been able to locate Proposition~\ref{ConsistencyIntro} anywhere in the lite\-ra\-ture, in an equivalent form a close variant of it, for the analogous measure on real symmetric positive-definite matrices, appears to be folklore in the statistical literature, see for example \cite[Chapter~3]{StatisticalMatrixAnalysis}. The argument there makes use of the technique of Schur complementation, the formula for the determinant of a block matrix and the special form of the density of $\mathcal{M}^{(\nu),N}$. The proof presented in Section~\ref{section2} is completely different and makes a novel use of the connection to orthogonal polynomials.
\end{Remark}

\subsection[Description of ergodic decomposition of $\mathsf{M}^{(\nu)}$]{Description of ergodic decomposition of $\boldsymbol{\mathsf{M}^{(\nu)}}$}\label{section1.4}

Consider the measure $\mathsf{M}^{(\nu)}$ defined above for $\nu >-1$. It is a result of Borodin and Olshanski, see Proposition~4.4 in~\cite{BorodinOlshanskiErgodic}, that there exists a unique probability measure $\mathfrak{m}^{(\nu)}$ on $\Omega$ such that
\begin{gather*}
\mathsf{M}^{(\nu)}=\int_{\Omega}^{}M_{\omega}\mathfrak{m}^{(\nu)}({\rm d}\omega).
\end{gather*}
Here, the equality is interpreted as integrated against a test function, see Section~\ref{section3}.
The main result of this paper is the explicit description of~$\mathfrak{m}^{(\nu)}$. To proceed and state it precisely we need some more definitions and background on determinantal measures.

\begin{Definition} Define the subset $\Omega_0^+ \subset \Omega$ such that $\omega=\left(\alpha^+,\alpha^-, \gamma_1, \delta\right) \in \Omega_0^+$ iff
\begin{gather*}
\alpha_i^{-}(\omega) \equiv 0, \qquad \alpha_j^{+}(\omega) \neq 0, \qquad \text{for all} \quad i,j \in \mathbb{N};\\
\gamma_2(\omega) =\delta(\omega)-\sum_{i}^{}\big(\alpha_i^+(\omega)\big)^2-\sum_{i}^{}\big(\alpha_i^-(\omega)\big)^2=0;\\
\gamma_1(\omega) = \sum_{i }^{}\alpha_i^{+}(\omega)<\infty.
\end{gather*}
\end{Definition}
In particular, on $\Omega_0^+$ there is no `Gaussian component' namely $\gamma_2(\omega)\equiv 0$ and the para\-me\-ter~$\gamma_1(\omega)$ is completely determined by the $\alpha_i^+(\omega)$.

\paragraph{Determinantal measures}
Let $\mathcal{X}$ be a locally compact Polish space equipped with a $\sigma$-finite reference measure $\mu$. Let $\mathsf{Conf}(\mathcal{X})$ be the space of point configurations over $\mathcal{X}$. Points in a point configuration $\mathfrak{X} \in \mathsf{Conf}(\mathcal{X})$ will be called particles. We can embed $\mathsf{Conf}(\mathcal{X})$ in the space of finite measures on $\mathcal{X}$ by $\mathfrak{X}\mapsto \sum\limits_{x \in \mathfrak{X}}^{}\delta_x$ and with the induced topology $\mathsf{Conf}(\mathcal{X})$ is a Polish space.

A Borel probability measure $\mathbb{P}$ on $\mathsf{Conf}(\mathcal{X})$ is called a determinantal point process or measure, see \cite{Soshnikov}, with (Hermitian) kernel $\mathsf{K}\colon \mathcal{X} \times \mathcal{X} \to \mathbb{C}$ if for any $n\in \mathbb{N}$ and function $\Phi\in C_c\big(\mathcal{X}^n\big)$, continuous and of compact support, we have
\begin{gather}
\int_{\mathsf{Conf}(\mathcal{X})}^{} \sum_{x_1, \dots,x_n \in \mathfrak{X}}\Phi(x_1,\dots,x_n)\mathbb{P}({\rm d}\mathfrak{X})\nonumber\\
\qquad{} =\int_{\mathcal{X}^n}^{}\Phi(x_1,\dots,x_n)\det (\mathsf{K}(x_i,x_j) )_{i,j=1}^n{\rm d}\mu(x_1)\cdots {\rm d}\mu(x_n),\label{determinantaldefinition}
\end{gather}
where the sum is taken over ordered $n$-tuples of particles with pairwise distinct labels. The measure $\mathbb{P}$ satisfying (\ref{determinantaldefinition}) completely determines the pair $(\mathsf{K},\mu)$ and dropping dependence on $\mu$ (usually fixed), we denote it by $\mathbb{P}_{\mathsf{K}}$.

Throughout this paper the reference measure will be Lebesgue measure. We now define a~distinguished determinantal measure on $(0,\infty)$.

\begin{Definition}Define the Bessel kernel $\mathbb{J}_{\nu}$ by
\begin{gather*}
\mathbb{J}_{\nu}(x,y)=\frac{\sqrt{x}J_{\nu+1}(\sqrt{x})J_{\nu}(\sqrt{y})-\sqrt{y}J_{\nu+1}(\sqrt{y})J_{\nu}(\sqrt{x})}{2(x-y)},
\end{gather*}
where $J_{\nu}$ is the Bessel function of order $\nu$. This kernel gives rise to the Bessel determinantal point process $\mathbb{P}_{\mathbb{J}_{\nu}}$ on $(0,\infty)$ of Tracy and Widom \cite{TracyWidom}.
\end{Definition}

We are ready to give the explicit description of the ergodic decomposition of the inverse Wishart measures $\mathsf{M}^{(\nu)}$.

\begin{Theorem}\label{MainTheorem}
Let $\nu> -1$. The spectral measure $\mathfrak{m}^{(\nu)}$ associated to $\mathsf{M}^{(\nu)}$ is concentrated on~$\Omega_0^+$
\begin{gather*}
\mathfrak{m}^{(\nu)}\big(\Omega_0^+\big)=1.
\end{gather*}
Moreover, the law of the parameters $\alpha^+=\big(\alpha_1^+\ge \alpha_2^+ \ge \alpha_3^+ \ge \cdots \ge 0 \big)$ under $\mathfrak{m}^{(\nu)}$ viewed as a~point configuration on $(0,\infty)$ is given by the determinantal point process $\mathbb{P}_{K^{\nu}_{\infty}}$ with correlation kernel $K_{\infty}^{\nu}(x,y)$ defined by
\begin{gather*}
K_{\infty}^{\nu}(x,y)=\frac{8}{xy} \mathbb{J}_{\nu}\left(\frac{8}{x},\frac{8}{y}\right).
\end{gather*}
\end{Theorem}

Theorem \ref{MainTheorem} above will be proven by an approximation procedure, the so-called `ergodic method' of Olshanski and Vershik \cite{OlshanskiVershik}, that is explained in Section~\ref{section3}. For this asymptotic analysis we will exploit the relation with the Laguerre ensemble.
